\documentclass[11pt]{article}
\usepackage[utf8]{inputenc}
\usepackage[margin=1in]{geometry}
\usepackage{amsmath}
\usepackage{graphicx}
\usepackage{booktabs}
\usepackage{hyperref}
\usepackage[round]{natbib}

\title{Activity-Dependent Blood--Brain Barrier Modulation Underlies Synaptic Plasticity in the Healthy Somatosensory Cortex}
\author{Author A,$^{1}$ Author B,$^{2}$ Author C,$^{1}$ Author D,$^{3}$ Author E$^{1,*}$\\
\small $^{1}$Department of Neuroscience, Ben-Gurion University\\
\small $^{2}$Department of Biomedical Engineering, Ben-Gurion University\\
\small $^{3}$Department of Neurology, University Hospital\\
\small $^{*}$Corresponding author: authore@bgu.edu}
\date{}

\begin{document}
\maketitle

\begin{abstract}
The blood--brain barrier (BBB) has traditionally been considered a static protective boundary that prevents serum proteins from entering the brain parenchyma. BBB dysfunction following brain injury is known to drive pathological plasticity and epileptogenesis through TGF-$\beta$-dependent albumin signaling. Here we demonstrate that physiological neuronal activity induces focal, transient BBB modulation in the healthy brain, and that this modulation plays a mechanistic role in synaptic plasticity. Prolonged limb stimulation (30 minutes, 6~Hz) in rats induced extravasation of sodium fluorescein, Evans blue, and Alexa488-conjugated albumin specifically in the contralateral somatosensory cortex. ELISA confirmed elevated albumin concentration peaking at 30~minutes post-stimulation and declining over 4--24~hours---at levels significantly below those observed in photothrombosis injury controls. The stimulation produced long-term potentiation (LTP) of somatosensory evoked potentials, and direct cortical perfusion of albumin alone was sufficient to induce LTP. Pharmacological blockade of AMPA/NMDA receptors (CNQX/AP5), caveolae-mediated transcytosis (methyl-$\beta$-cyclodextrin), and TGF-$\beta$ receptor signaling (SJN) each abolished both BBB modulation and LTP, without disrupting hemodynamic responses (for mBCD and SJN). RNA-sequencing revealed upregulation of caveolae-associated transport and synaptic plasticity genes. Critically, we provide the first evidence in humans that 30~minutes of motor activity increases BBB permeability in the activated cortical region, as demonstrated by dynamic contrast-enhanced MRI co-localized with fMRI activation. These findings establish activity-dependent BBB modulation as a physiological mechanism underlying synaptic plasticity in the healthy brain.
\end{abstract}

\section{Introduction}

The blood--brain barrier (BBB) is a highly selective interface between the cerebral vasculature and the brain parenchyma, formed by specialized endothelial cells with tight junctions, surrounded by pericytes and astrocytic end-feet \citep{abbott2010structure, zlokovic2008blood}. Under normal conditions, the BBB severely restricts the passage of serum proteins, immune cells, and other circulating molecules into the brain, thereby maintaining the tightly controlled ionic and molecular environment essential for synaptic transmission and neuronal function \citep{pardridge2005molecular}.

BBB disruption has been extensively studied in pathological contexts. Following brain injury, stroke, or seizures, BBB breakdown permits the extravasation of serum albumin into the neuropil, where it activates astrocytic TGF-$\beta$ receptors, triggering a cascade of transcriptional changes that lead to hyperexcitability and epileptogenesis \citep{friedman2009blood, ivens2007tgf}. This pathological BBB-to-plasticity pathway has been well characterized and represents an important mechanism of disease progression. However, whether physiological levels of neuronal activity can modulate BBB permeability in the healthy, uninjured brain has remained an open question.

Several lines of evidence have hinted at physiological BBB modulation. Neurovascular coupling---the process by which neuronal activity regulates local cerebral blood flow---involves intimate signaling between neurons, astrocytes, and the vasculature \citep{iadecola2017neurovascular, attwell2010glial}. Pericytes, which ensheath brain capillaries, have been shown to regulate BBB permeability through modulation of transcytosis and tight junction integrity \citep{armulik2010pericytes}. Furthermore, caveolae-mediated transcytosis has emerged as a regulated pathway for transport across the BBB endothelium \citep{andreone2017blood, knowland2014stepwise}. Despite these advances, a direct demonstration that physiological neuronal activity induces BBB modulation and that this modulation contributes to normal synaptic plasticity has been lacking.

Long-term potentiation (LTP), the persistent strengthening of synaptic transmission following patterned activity, is the best-characterized cellular mechanism of learning and memory \citep{bliss1973long}. While canonical LTP mechanisms involve glutamate receptor activation, calcium influx, and postsynaptic signaling cascades, accumulating evidence suggests that the extracellular environment---including the availability of serum-derived factors---may play a role in plasticity. TGF-$\beta$ signaling, which is activated by albumin extravasation in pathological contexts, has also been implicated in physiological synaptic plasticity \citep{toth2019tgfbeta}. This raises the tantalizing possibility that regulated, physiological BBB modulation could serve as a normal plasticity mechanism, with pathological BBB breakdown representing an extreme version of this process.

In this study, we test the hypothesis that physiological neuronal activity induces focal BBB modulation in the healthy somatosensory cortex of rats, and that this BBB modulation contributes to LTP through caveolae-mediated albumin transcytosis and TGF-$\beta$ receptor signaling. We further test whether this phenomenon can be detected in humans using non-invasive MRI. Our results establish activity-dependent BBB modulation as a previously unrecognized physiological mechanism that links sustained neuronal activity to synaptic plasticity.

\section{Methods}

\subsection{Animal Preparation}

Adult male Sprague Dawley rats (300--350~g; Harlan Laboratories) were anesthetized with ketamine (100~mg/mL, 0.08~mL/100~g) and xylazine (20~mg/mL, 0.06~mL/100~g) intraperitoneally. Body temperature was maintained at $37 \pm 0.5$\textdegree C with a feedback-controlled heating pad. Heart rate, breath rate, and O$_2$ saturation were continuously monitored using a MouseOx pulse oximeter. A craniotomy was performed over the right sensorimotor cortex (4~mm caudal, 2~mm frontal, 5~mm lateral to bregma), and artificial cerebrospinal fluid (aCSF) was continuously applied to the exposed cortex.

\begin{table}[ht]
\centering
\caption{Artificial CSF Composition}
\label{tab:acsf}
\begin{tabular}{lc}
\toprule
\textbf{Component} & \textbf{Concentration (mM)} \\
\midrule
NaCl & 124 \\
NaHCO$_3$ & 26 \\
NaH$_2$PO$_4$ & 1.25 \\
MgSO$_4$ & 2 \\
CaCl$_2$ & 2 \\
KCl & 3 \\
Glucose & 10 \\
\bottomrule
\end{tabular}
\end{table}

\subsection{Stimulation Protocol}

Limb stimulation was delivered via two subdermal needle electrodes inserted into the left forelimb or hindlimb. Stimulation consisted of 0.1~ms square pulses at 2--3~mA and 6~Hz. Test stimulations (360 pulses, 60~s) were delivered before and after a 30-minute conditioning stimulation at 6~Hz.

\subsection{BBB Permeability Assessment}

Three complementary approaches assessed BBB permeability. First, sodium fluorescein (NaFlu; MW $\sim$376~Da) was injected intravenously and detected by intravital fluorescence microscopy (wide-field and two-photon) during and after stimulation. Second, Evans blue (EB; binds albumin to form a $\sim$67~kDa complex) and Alexa Fluor 488-conjugated albumin ($\sim$67~kDa) were injected intravenously in separate cohorts without craniotomy, and extravasation was assessed histologically. Third, cortical tissue was harvested at 30~minutes, 4~hours, and 24~hours post-stimulation for ELISA quantification of albumin concentration. Photothrombosis-induced ischemia (24~hours post-injury) served as a positive control for the ELISA assay.

\subsection{Electrophysiology}

Somatosensory evoked potentials (SEPs) were recorded as local field potentials from layer 2/3 of the sensorimotor cortex. SEP amplitude was measured in response to 1-minute test stimulations (6~Hz, 2~mA) delivered before and after the 30-minute conditioning stimulation. LTP was defined as a significant increase in SEP amplitude post-conditioning relative to baseline.

\subsection{Pharmacological Experiments}

Three pharmacological blockers were applied via cortical perfusion: CNQX/AP5 (50~$\mu$M; AMPA/NMDA glutamate receptor antagonists), methyl-$\beta$-cyclodextrin (mBCD; 10~$\mu$M; caveolae/lipid raft disruptor), and SJN (0.3~mM; TGF-$\beta$ receptor inhibitor). Each blocker was applied before and during the 30-minute stimulation, and effects on BBB permeability (NaFlu extravasation), synaptic plasticity (SEP amplitude), and hemodynamic response (total hemoglobin, HbT) were assessed.

\subsection{RNA-Sequencing}

Contralateral somatosensory cortex tissue was collected at 1~hour and 24~hours post-stimulation ($n = 8$ rats per group). Differential gene expression analysis compared 24-hour versus 1-hour timepoints using the Wald test. Gene Ontology (GO) enrichment analysis was performed on significant differentially expressed genes.

\subsection{Human Experiments}

Human subjects performed a 30-minute stress-ball squeezing task with one hand. Functional MRI (fMRI) localized the activated cortical region, and dynamic contrast-enhanced (DCE) MRI assessed BBB permeability before and after the motor task. Permeability maps were compared within the fMRI-defined activation region.

\section{Results}

\subsection{Prolonged Stimulation Induces Focal BBB Permeability in the Contralateral Cortex}

Thirty minutes of limb stimulation at 6~Hz induced focal BBB permeability increase in the contralateral somatosensory cortex. Intravital fluorescence microscopy revealed NaFlu extravasation around responding blood vessels specifically in the stimulated cortical region, while the ipsilateral cortex showed no tracer leakage. Control animals (no 30-minute stimulation) showed no extravasation. In separate cohorts without craniotomy, both Evans blue and Alexa488-conjugated albumin were detected in the neuropil of the contralateral hemisphere, confirming that serum albumin crosses the BBB following sustained physiological stimulation (Table~\ref{tab:tracers}).

\begin{table}[ht]
\centering
\caption{BBB Tracer Extravasation Following 30-Minute Stimulation}
\label{tab:tracers}
\begin{tabular}{llcc}
\toprule
\textbf{Tracer} & \textbf{MW Class} & \textbf{Contralateral} & \textbf{Ipsilateral/Control} \\
\midrule
Sodium fluorescein & Low ($\sim$376 Da) & Extravasation & No extravasation \\
Evans blue & High ($\sim$67 kDa complex) & Extravasation & No extravasation \\
Alexa488-albumin & High ($\sim$67 kDa) & Extravasation & No extravasation \\
\bottomrule
\end{tabular}
\end{table}

ELISA quantification showed that albumin concentration in stimulated cortex was significantly elevated at 30~minutes post-stimulation compared to sham controls, declined by 4~hours, and returned to near-baseline by 24~hours. Importantly, albumin levels at all time points were significantly lower than those observed in photothrombosis injury controls, demonstrating that the activity-dependent BBB modulation is transient and far milder than pathological BBB breakdown.

\subsection{Stimulation-Induced BBB Modulation Accompanies LTP}

The 30-minute conditioning stimulation produced LTP of somatosensory evoked potentials, manifested as a significant increase in SEP amplitude when assessed with 1-minute test stimulations delivered after the conditioning period. Sham controls showed no change in SEP amplitude. Strikingly, direct cortical perfusion of albumin (without any peripheral stimulation) was also sufficient to induce LTP, replicating the potentiation effect and providing direct evidence that albumin extravasation mediates the plasticity response.

\subsection{Pharmacological Dissection of the BBB--Plasticity Pathway}

To identify the molecular mechanisms linking neuronal activity, BBB modulation, and LTP, we tested three pharmacological blockers applied via cortical perfusion (Table~\ref{tab:blockers}).

\begin{table}[ht]
\centering
\caption{Effects of Pharmacological Blockers on BBB Permeability, LTP, and Hemodynamics}
\label{tab:blockers}
\begin{tabular}{lcccc}
\toprule
\textbf{Condition} & \textbf{BBB Permeability} & \textbf{LTP} & \textbf{Hemodynamic} \\
\midrule
Stim alone & Increased & Present & Present \\
CNQX/AP5 + stim & Blocked & Blocked & Blocked \\
mBCD + stim & Blocked & Blocked & Not affected \\
SJN + stim & Blocked & Blocked & Not affected \\
Albumin perfusion & N/A & Present & N/A \\
\bottomrule
\end{tabular}
\end{table}

CNQX/AP5 (glutamate receptor antagonists) abolished both BBB modulation and LTP, confirming that excitatory neurotransmission is required for the entire cascade. mBCD (caveolae disruptor) blocked both BBB opening and LTP without significantly affecting the hemodynamic response (HbT signal), indicating that caveolae-mediated transcytosis is required for activity-dependent BBB modulation and that this pathway is separable from vasodilatory neurovascular coupling. SJN (TGF-$\beta$ receptor inhibitor) similarly blocked both BBB modulation and LTP without disrupting hemodynamics, implicating TGF-$\beta$ signaling as a downstream mediator of albumin-induced plasticity.

The critical finding that mBCD and SJN block BBB modulation and LTP without disrupting the hemodynamic response demonstrates that the BBB--plasticity pathway is mechanistically distinct from the vasodilatory component of neurovascular coupling.

\subsection{Transcriptomic Evidence for BBB Transport and Plasticity Gene Upregulation}

RNA-sequencing of stimulated cortical tissue revealed significant differentially expressed genes (DEGs) between 24-hour and 1-hour post-stimulation timepoints ($n = 8$ per group, Wald test). Gene Ontology enrichment analysis of upregulated genes at 24~hours identified enrichment for BBB transport-related terms (vesicle-mediated transport, transcytosis, including caveolae-associated genes) and synaptic plasticity-related terms (postsynaptic density, synapse organization, long-term potentiation). At the 1-hour timepoint, immediate early genes and activity-dependent transcription factors were transiently upregulated, consistent with an initial neuronal activity response. The convergence of BBB transport and synaptic plasticity ontologies in the transcriptomic data provides molecular-level support for the mechanistic link between activity-dependent BBB modulation and LTP.

\subsection{Human Motor Activity Increases BBB Permeability in the Activated Cortex}

To determine whether activity-dependent BBB modulation occurs in humans, subjects performed a 30-minute stress-ball squeezing task. fMRI confirmed activation of the contralateral motor and somatosensory cortex during the task. DCE-MRI permeability maps acquired before and after the task showed increased BBB permeability in the fMRI-defined activation region following the 30-minute motor task. The spatial co-localization of increased BBB permeability with the fMRI activation map provides the first evidence that physiological motor activity modulates BBB permeability in the human brain.

\subsection{Proposed Mechanistic Pathway}

Integrating across all experimental modalities, we propose the following mechanistic pathway (Table~\ref{tab:pathway}):

\begin{table}[ht]
\centering
\caption{Proposed Mechanistic Pathway for Activity-Dependent BBB-Mediated Plasticity}
\label{tab:pathway}
\begin{tabular}{clp{7cm}}
\toprule
\textbf{Step} & \textbf{Component} & \textbf{Evidence} \\
\midrule
1 & Sustained neuronal activity & LFP recordings confirm cortical activation \\
2 & Glutamate release, AMPA/NMDA activation & CNQX/AP5 blocks both BBB opening and LTP \\
3 & Caveolae-mediated albumin transcytosis & mBCD blocks BBB opening and LTP \\
4 & Albumin enters neuropil & ELISA, EB, Alexa488-albumin confirm extravasation \\
5 & TGF-$\beta$ receptor signaling & SJN blocks BBB opening and LTP \\
6 & Synaptic potentiation (LTP) & SEP increase; albumin perfusion sufficient \\
7 & Gene expression changes & RNA-seq: BBB transport and plasticity genes upregulated \\
\bottomrule
\end{tabular}
\end{table}

\section{Discussion}

Our findings challenge the traditional view of the blood--brain barrier as a static protective boundary and establish activity-dependent BBB modulation as a physiological mechanism contributing to synaptic plasticity in the healthy brain. This work represents a conceptual shift in understanding the relationship between the BBB and brain function.

The demonstration that 30 minutes of sustained limb stimulation induces focal, transient albumin extravasation in the contralateral somatosensory cortex provides the first direct evidence of physiological BBB modulation in the healthy brain. Several features distinguish this physiological modulation from pathological BBB breakdown. First, the albumin levels are significantly lower than those observed in photothrombosis injury, indicating a quantitatively different process. Second, the modulation is transient, peaking at 30~minutes and declining over hours, rather than persisting or progressing as in injury. Third, the modulation is focal, confined to the activated cortical region, and lateralized to the contralateral hemisphere.

The pharmacological experiments reveal a specific molecular cascade: neuronal activity drives glutamate release that activates AMPA/NMDA receptors, triggering caveolae-mediated transcytosis of serum albumin across the BBB endothelium, followed by TGF-$\beta$ receptor activation and downstream plasticity. The finding that mBCD and SJN block both BBB modulation and LTP without disrupting the hemodynamic response demonstrates that the BBB--plasticity pathway is mechanistically separable from the vasodilatory component of neurovascular coupling \citep{iadecola2017neurovascular, attwell2010glial}. This is an important distinction, as it suggests that neurovascular coupling encompasses at least two distinct functional arms: regulation of blood flow (vasoregulatory) and regulation of BBB permeability (transendothelial transport).

The transcriptomic data provide molecular-level confirmation: genes associated with caveolae-mediated transport and synaptic plasticity are upregulated following stimulation, supporting the proposed mechanistic pathway. The temporal progression from immediate early gene activation at 1~hour to BBB transport and plasticity gene upregulation at 24~hours is consistent with a cascade in which acute neuronal activity triggers transcriptional programs that sustain and consolidate the plasticity response.

Perhaps the most compelling evidence comes from the human DCE-MRI data, which demonstrate that 30~minutes of hand squeezing a stress ball increases BBB permeability in the activated motor cortex region. The spatial co-localization with fMRI activation provides strong evidence that activity-dependent BBB modulation is a conserved physiological mechanism in the human brain.

These findings have several important implications. First, they suggest that the pathological BBB--TGF-$\beta$ signaling cascade previously described in epileptogenesis \citep{friedman2009blood, ivens2007tgf} may represent a dysregulated extreme of a normally beneficial physiological process. Second, they identify the BBB as an active participant in synaptic plasticity, adding a vascular dimension to the cellular mechanisms of learning and memory \citep{bliss1973long}. Third, they raise the possibility that disruption of activity-dependent BBB modulation could contribute to cognitive decline in conditions affecting the cerebrovasculature.

Limitations include the use of anesthesia in rat experiments, which may affect the magnitude of BBB modulation, and the relatively small sample size for the human MRI study. Future work should investigate whether learning tasks produce similar BBB modulation and whether this mechanism is affected by aging or cerebrovascular disease.

\section{Conclusion}

This study demonstrates that physiological neuronal activity induces focal, transient modulation of blood--brain barrier permeability in the healthy somatosensory cortex, and that this modulation mediates synaptic long-term potentiation through a pathway involving caveolae-mediated albumin transcytosis and TGF-$\beta$ receptor signaling. The pathway is mechanistically separable from hemodynamic neurovascular coupling, as pharmacological disruption of caveolae or TGF-$\beta$ receptors blocks both BBB modulation and LTP without affecting vasodilation. RNA-sequencing confirms upregulation of BBB transport and plasticity genes, and DCE-MRI provides the first evidence that this phenomenon occurs in the human brain during sustained motor activity. These findings establish activity-dependent BBB modulation as a physiological mechanism of synaptic plasticity, reconceptualizing the BBB as an active participant in brain function rather than a passive protective barrier.

\bibliographystyle{plainnat}
\bibliography{references}

\end{document}
